\section{Study Design}
\label{sec:design}

The study combines repository mining (selecting and labelling real optimization commits), benchmarking experiments (cost passes, an A/B slowdown oracle and an A/A control) and a quantitative simulation of reduction arms over the replayed edits.
We mine real commits rather than generating candidates with an optimizer because real commits come with an expert's diff and a stated intent, so whether an edit is an optimization can be judged, and no optimizer's biases enter the edit set; the mining rule's own biases are discussed in Section~\ref{sec:lessons}.
Fig.~\ref{fig:pipeline} shows the pipeline.
A written protocol (subjects, mining rule, arms, frozen fault sets, decision rules and timing protocol) was committed before the measurements \TODO{commit hash and date}.
Parts of the analysis are post hoc: the commit labels that define the RQ2 edit set were added after mining and after the first coverage count, the noise-run labelling rule was written after a first look at two slow states, and the share analysis and the profile re-check were added after the first results; all deviations are listed in Section~\ref{sec:deviations}.
The selection and validation rules are deterministic and model-free: no learned or LLM component decides what to build, test or benchmark. The commit labels were produced by LLM agents (Section~\ref{sec:threats}).

\begin{figure}[t]
\centering
\begin{tikzpicture}[node distance=3mm and 4mm, every node/.style={font=\footnotesize},
  box/.style={draw, rounded corners=2pt, align=center, minimum height=7mm, inner sep=2pt, text width=13mm},
  arr/.style={-{Stealth[length=4pt]}}]
\node[box] (mine) {mine\\(subject rule)};
\node[box, right=of mine] (label) {label\\(two raters)};
\node[box, right=of label] (replay) {replay\\edit sequence};
\node[box, right=of replay] (prof) {profile and\\static reach};
\node[box, below=6mm of replay, text width=20mm] (measure) {measure per state:\\build, tests, benchmarks};
\node[box, below=6mm of prof, text width=15mm] (inject) {frozen mutants,\\slowdowns, A/A};
\node[box, below=6mm of label, text width=17mm] (arms) {simulate arms\\(selection counts)};
\draw[arr] (mine) -- (label); \draw[arr] (label) -- (replay); \draw[arr] (replay) -- (prof);
\draw[arr] (replay) -- (measure); \draw[arr] (prof) -- (inject); \draw[arr] (measure) -- (arms);
\draw[arr] (inject) -- (measure);
\end{tikzpicture}
\caption{Study pipeline. Mining, labelling and replay fix the edit sequence; profiles and static reach give benchmark coverage (RQ2); measured states and injected faults give cost, detection and power (RQ1, RQ3, RQ4).}
\label{fig:pipeline}
\end{figure}

\subsection{Subjects}
\textbf{networkx} is the main subject, chosen because it has an asv suite in its own repository, a pure-Python build (so build cost does not mask test and benchmark cost) and a recent history with many commits that describe themselves as performance changes. It has an asv suite of \NxBenchCases{} benchmark cases in \NxBenchFiles{} files and \NxNumTests{} tests.
\textbf{bottleneck}, a NumPy extension written in C from templates, is a contrast subject for build cost only, on \BnEdits{} replayed edits (states after an uncompilable replayed state are held out, Section~\ref{sec:lessons}).
\TODO{Table 2: subjects and edit counts. Add the E1 replication projects (3--4, chosen by the pre-registered order and attrition rule; candidates xarray, scikit-image, astropy, geopandas, shapely, pandas) with mined, labelled, true-optimization and profiled counts.}

\subsection{Edit mining, labelling and replay}
\textbf{Mining.} A commit qualifies if its subject matches a fixed regular expression for performance wording (\texttt{perf}, \texttt{speed up}, \texttt{faster}, \texttt{optimize}, \texttt{performance}, \texttt{vectorize}, \texttt{efficien}) and it changes package source.
For networkx we took the latest \NxMined{} qualifying commits before the study started (November 2026), in chronological order, without manual inclusion or exclusion; the number was fixed in advance to bound the cost of full validation at every state.
\TODO{Give the date span of the mined commits and their size (changed lines and functions per edit) in Table 2.}

\textbf{Unit of analysis and changed functions.} The unit is one replayed edit (one commit).
A Python function counts as changed if any removed or added line of the diff falls in its span, including its decorators and its qualified name inside classes and nested functions; added and deleted functions count, blank and comment-only lines do not.
Changed lines outside every function are module-level changes; they are excluded from coverage counts and expanded to all functions of the module by the selectors.
For bottleneck, C function definitions are found by a brace scan and template headers map to the Python-visible name.
\TODO{Validate the changed-function extraction against a hand check of a sample of edits and report the agreement.}

\textbf{Labelling.} Keyword mining is imprecise (Section~\ref{sec:lessons}), so every mined commit of four projects (\LabelRows{} rows, including the Q3 pilot projects of Appendix~\ref{app:q3}) was labelled from its subject, body and diff into three categories (optimization, not an optimization, unsure), using a frozen codebook whose question is whether the commit's main intent is to make existing code faster or use fewer resources without changing what it computes.
Two raters labelled all rows blind to benchmark reach, replay status and each other; a third agent, the adjudicator, decided the \LabelAdjudicated{} disagreements.
The raters were two LLM agents of the same model family; we treat this as a threat (Section~\ref{sec:threats}).
No human rater labelled the rows; the labels are agent labels throughout.
\TODO{Name the model and version of the raters and the adjudicator, and put the prompts in the artifact.}
Cohen's $\kappa$ over the three categories was \LabelKappa{} (\LabelAgree{} of \LabelRows{} rows agree).
The primary edit set for RQ2 is the commits both raters or the adjudicator labelled as optimizations; an earlier reading of the commit subjects alone is reported as a sensitivity set.
\TODO{Replication (E1): codebook v2 frozen before mining, agent-only labels as here, and held-out precision of the v2 mining rule.}

\textbf{Replay.} The mined commits are applied in order on top of the parent of the oldest one, giving a sequence of states $S_0, S_1, \dots$; \NxDropped{} of the \NxMined{} networkx commits did not apply and were dropped by rule, leaving \NxReplayed{} replayed edits (\NxStates{} states).
Of these, \NxAgentOpt{} are labelled optimizations by the agents and \NxManualOpt{} by the subject-only reading.
The replayed sequence is not a real optimizer session: consecutive edits are historical commits, not minutes-apart candidates of one agent.
Results that are per edit (RQ2, RQ3) do not depend on this; cost totals and the accumulation policy (RQ4) do.

\subsection{Arms}
Every arm validates every replayed edit.
\textbf{FULL} (A0) rebuilds from clean, runs all tests serially and all benchmarks.
\textbf{INC} (A1) builds incrementally.
\textbf{TESTMON} (A2) runs \texttt{pytest --testmon} unchanged, with its database carried across steps, as PerfAgent does~\cite{deng2026perfagent}.
\textbf{CHANGED-MODULE} (A3) runs the tests whose file name maps to a changed module.
\textbf{STATIC-BENCH} (A5) selects benchmarks whose static call-graph reach intersects the changed functions.
\textbf{A4} is our deterministic policy: function-level test selection over the changes accumulated since the last green full check, profile-based benchmark selection, and a full check triggered when an additive change score crosses a fixed threshold (a model-free variant of risk-mass batching~\cite{sayedsalehi2026riskaware}); its parameters were fixed before the runs.
Arms are also combined with periodic full checks every $k$ edits.
Test counts per arm come from real selections (testmon) or from the measured dependency map; seconds are estimated from per-test durations at $S_0$ and calibrated against measured end-to-end time.

\subsection{Coverage measures}
A benchmark case \emph{executes} an edit if its measured profile contains at least one function the edit changed (module-level lines excluded).
Profiles were recorded with \texttt{sys.monitoring} for all \NxBenchCases{} cases at $S_0$, during one call of each benchmark after its setup, so functions that run only in untimed setup or at import do not count as executed; this is a counting run, not a timing run.
\emph{Static reach} uses a conservative import and call graph from $S_0$.
An edit is \emph{unreached} if no case executes it.
Reach was taken at $S_0$ for all edits, not at each edit's parent.
\TODO{E1: profiles at each edit's parent for the replication projects, with the grouped profile runner validated against the networkx $S_0$ reach sets.}

\subsection{Oracles}
\textbf{Functional.} \MutantsFrozen{} mutants were frozen before any arm was configured, each placed in a function changed at its step; \MutantsKilled{} are killed by the full suite and form the functional fault set.
\textbf{Performance.} \NumSlowdowns{} synthetic slowdowns were frozen, each in a function changed at its step: nine call the function body twice and one replaces a set by a list for membership tests.
They are probes of detectability, not models of real regressions.
For each slowdown, a patched and an unpatched clone are measured interleaved (A/B), and a case is \emph{slower} under asv's rule: ratio of medians above \AsvFactor{} and a two-sided Mann-Whitney $p$ below \AsvAlpha~\cite{asvcommands}.
All primary verdicts use the exact Mann-Whitney test; re-testing all \RecomputeVerdicts{} verdicts with the exact test changed none of them (Section~\ref{sec:rq3}).
Within reach, each side has two processes with three values each, in three interleaved repeats; the profile-completeness check and the A/A control use five processes with three values per side and repeat.
A pre-registered noise floor additionally requires the ratio to exceed $1+3\sigma$, with $\sigma$ from a same-code A/A control.
The protocol also fixed strength bands for the regression set used to judge selection safety: a ``strong'' slowdown makes some case at least twice as slow, a ``moderate'' one between \BandModLo{} and \BandModHi\x.
\textbf{Test assumptions.} Mann-Whitney assumes independent samples; samples from one process are correlated, which makes $p$ optimistic, and repeating the oracle (Section~\ref{sec:rq3}) is the check. Wilson intervals assume independent trials; edits and slowdowns that share functions are not fully independent, so the intervals are optimistic.
\textbf{A/A control.} A frozen, seeded sample of \AACases{} cases (\AAAlgoCases{} of them from the heavy module) was run as $S_0$ against $S_0$ in two clones, with both run orders (ABBA), so every non-``same'' verdict is false.
Proportions are reported with Wilson 95\pct{} intervals.
\textbf{Exact test.} The rule we first implemented used scipy's default Mann-Whitney method, a normal approximation at these sample sizes, while asv uses an exact test; re-testing all \RecomputeVerdicts{} verdicts from their raw samples under the exact test changed \RecomputeChanged{} of them, and borderline $p$ values became smaller. All $p$ values in this paper are exact.
The smallest attainable two-sided exact $p$ is \MinPSixSix{} for six against six samples, above $\alpha=\AsvAlpha$ (Section~\ref{sec:lessons}), and \MinPFifteen{} for fifteen against fifteen.

\subsection{Platform and timing protocol}
All runs used one desktop machine (\texttt{Intel i5-12400F}, \TODO{core/thread and memory counts as macros}), WSL2 \texttt{Ubuntu 24.04}, \texttt{CPython 3.12}, with the package versions frozen in the artifact.
The CPU governor cannot be pinned under WSL2.
Timing runs start only when the machine is idle by a load gate (one-minute load average at most \LoadGate), and runs whose load exceeds the gate are labelled noisy and reported as such, not dropped.
Benchmarks use our asv-compatible runner with asv's default protocol for cost passes and more processes for verdicts.
Counts (reach, selections, detections) are machine-independent; wall times are reported as ranges over passes or as a median with its outliers named, never as single-pass point values.

\textbf{Ethics.} We analyse public open-source history and our own measurements; no personal data are analysed and no developers are named or assessed.
